\section{Functorial resolutions and finite filtered injectives}\label{proofs-7589-8350-functorial-resolutions-and-finite-filtered-injectives}

Editorial review of the Stacks Project authors\textquotesingle{} Derived Categories,\allowbreak{} original commit a04446e5\allowbreak{}7ec1fbc2\allowbreak{}52a871af\allowbreak{}cec7752f\allowbreak{}b2807b14, lines 7589–\allowbreak{}8350. The sources and actual reading ranges are bound in SOURCE\_\allowbreak{}READING\_\allowbreak{}PART11.json. The original formulas and objects remain the comparison authority. The stronger statements below are editorial consequences,\allowbreak{} with no novelty claim or replacement of source exposition. Earlier complete proofs are retained in PROOFS\_\allowbreak{}6209\_\allowbreak{}7089.md,\allowbreak{} PROOFS\_\allowbreak{}7096\_\allowbreak{}7578.md,\allowbreak{} and the earlier bounded-filtered-category review.

\subsection{\texorpdfstring{1. The embedding functor,\allowbreak{} its zero summand,\allowbreak{} and every differential}{1. The embedding  its zero  and every differential}}\label{proofs-7589-8350-the-embedding-functor-its-zero-summand-and-every-differential}

Current Homology 7232–\allowbreak{}7248 defines a functorial injective embedding as a functor to the arrow category. Equivalently it gives an ordinary functor J:\allowbreak{}A\allowbreak{}->A and natural monomorphisms eta\_\allowbreak{}A:\allowbreak{}A\allowbreak{}->J(A)\allowbreak{},\allowbreak{} with each J(A)\allowbreak{} injective. It does not require J to be additive.

Keep the original J and its zero-object contribution. Write s\_\allowbreak{}A\allowbreak{}=J(0\allowbreak{}->A)\allowbreak{}:\allowbreak{}J(0)\allowbreak{}\allowbreak{}->J(A)\allowbreak{} and p\_\allowbreak{}A\allowbreak{}=J(A\allowbreak{}->0)\allowbreak{}:\allowbreak{}J(A)\allowbreak{}\allowbreak{}->J(0)\allowbreak{}. Their composite p\_\allowbreak{}A s\_\allowbreak{}A is 1. Put J\_\allowbreak{}red(A)\allowbreak{}\allowbreak{}=ker p\_\allowbreak{}A,\allowbreak{} with inclusion k\_\allowbreak{}A. There is a unique r\_\allowbreak{}A:\allowbreak{}J(A)\allowbreak{}\allowbreak{}->J\_\allowbreak{}red(A)\allowbreak{} satisfying k\_\allowbreak{}A r\_\allowbreak{}A\allowbreak{}=1-s\_\allowbreak{}Ap\_\allowbreak{}A. The two inverse isomorphisms are

\[
 (s_A,k_A):J(0)\oplus J_{\rm red}(A)\longrightarrow J(A),
 \qquad
 \binom{p_A}{r_A}:J(A)\longrightarrow J(0)\oplus J_{\rm red}(A).
                                                               \tag{1.1}
\]

For f:\allowbreak{}A\allowbreak{}->B,\allowbreak{} p\_\allowbreak{}B J(f)\allowbreak{}\allowbreak{}=p\_\allowbreak{}A and J(f)\allowbreak{}s\_\allowbreak{}A\allowbreak{}=s\_\allowbreak{}B. Thus J(f)\allowbreak{} restricts to a map J\_\allowbreak{}red(f)\allowbreak{},\allowbreak{} and (1.1)\allowbreak{} identifies it with the matrix diag(1\_\allowbreak{}\{J(0)\allowbreak{}\},\allowbreak{}J\_\allowbreak{}red(f)\allowbreak{})\allowbreak{}. These restrictions preserve identity and composition. J\_\allowbreak{}red(0)\allowbreak{}\allowbreak{}=ker(1\_\allowbreak{}\{J(0)\allowbreak{}\})\allowbreak{}\allowbreak{}=0. Its values are injective direct summands:\allowbreak{} extend a map into J\_\allowbreak{}red(A)\allowbreak{} into J(A)\allowbreak{} and then compose with r\_\allowbreak{}A.

Naturality gives p\_\allowbreak{}A eta\_\allowbreak{}A\allowbreak{}=eta\_\allowbreak{}0(A\allowbreak{}->0)\allowbreak{}\allowbreak{}=0. Therefore eta\_\allowbreak{}A factors uniquely as k\_\allowbreak{}A eta\_\allowbreak{}red,\allowbreak{}A; the factor is monic. The factorization and its naturality are retained. We use J\_\allowbreak{}red for the source\textquotesingle s explicitly described replacement,\allowbreak{} without claiming that the original J(0)\allowbreak{} was absent. Because a zero morphism factors through 0,\allowbreak{} J\_\allowbreak{}red sends every zero morphism to zero. This is the property needed for the following complexes,\allowbreak{} rather than additivity.

Define functors Q\_\allowbreak{}0(A)\allowbreak{}\allowbreak{}=A and recursively \[
 E_q(A)=J_{\rm red}(Q_q(A)),\qquad
 Q_{q+1}(A)=\operatorname{coker}
       (\eta_{\rm red,Q_q(A)}:Q_q(A)\to E_q(A)).
                                                               \tag{1.2}
\] Let pi\_\allowbreak{}q:\allowbreak{}E\_\allowbreak{}q\allowbreak{}->Q\_\allowbreak{}\{q\allowbreak{}+1\} be the quotient and v\_\allowbreak{}q\allowbreak{}=eta\_\allowbreak{}red,\allowbreak{}Q\_\allowbreak{}\{q\allowbreak{}+1\} pi\_\allowbreak{}q:\allowbreak{}E\_\allowbreak{}q\allowbreak{}->E\_\allowbreak{}\{q\allowbreak{}+1\}. Cokernel maps are induced by the natural commutative squares,\allowbreak{} so Q\_\allowbreak{}q,\allowbreak{} E\_\allowbreak{}q,\allowbreak{} pi\_\allowbreak{}q,\allowbreak{} and v\_\allowbreak{}q are functorial. They preserve zero objects and zero maps. We have v\_\allowbreak{}\{q\allowbreak{}+1\}v\_\allowbreak{}q\allowbreak{}=0,\allowbreak{} ker v\_\allowbreak{}q\allowbreak{}=im eta\_\allowbreak{}red,\allowbreak{}Q\_\allowbreak{}q,\allowbreak{} and im v\_\allowbreak{}\{q-1\}\allowbreak{}=im eta\_\allowbreak{}red,\allowbreak{}Q\_\allowbreak{}q for q>\allowbreak{}=1. The initial kernel is the image of A. Thus \[
 0\to A\xrightarrow{\eta_{\rm red,A}}E_0(A)
     \xrightarrow{v_0}E_1(A)\xrightarrow{v_1}\cdots             \tag{1.3}
\] is an injective resolution,\allowbreak{} naturally in A.

For the original bounded-below K,\allowbreak{} with K\textasciicircum{}p\allowbreak{}=0 for p<N,\allowbreak{} take I\textasciicircum{}\{p,\allowbreak{}q\}\allowbreak{}=E\_\allowbreak{}q(K\textasciicircum{}p)\allowbreak{} for q>\allowbreak{}=0 and zero otherwise,\allowbreak{} d\_\allowbreak{}1\textasciicircum{}\{p,\allowbreak{}q\}\allowbreak{}=E\_\allowbreak{}q(d\_\allowbreak{}K\textasciicircum{}p)\allowbreak{},\allowbreak{} and d\_\allowbreak{}2\textasciicircum{}\{p,\allowbreak{}q\}\allowbreak{}=v\_\allowbreak{}q(K\textasciicircum{}p)\allowbreak{}. Functoriality and preservation of zero prove d\_\allowbreak{}1\textasciicircum{}2\allowbreak{}=0. Equation (1.3)\allowbreak{} proves d\_\allowbreak{}2\textasciicircum{}2\allowbreak{}=0; naturality of v\_\allowbreak{}q proves d\_\allowbreak{}1d\_\allowbreak{}2\allowbreak{}=d\_\allowbreak{}2d\_\allowbreak{}1. All columns p\textless N vanish. The augmentation epsilon\textasciicircum{}p\allowbreak{}=eta\_\allowbreak{}red,\allowbreak{}K\textasciicircum{}p is a horizontal chain map and is killed by d\_\allowbreak{}2. The total complex has the exact differential \[
 d_{\operatorname{Tot}}|_{I^{p,q}}=d_1^{p,q}+(-1)^p d_2^{p,q}.
                                                               \tag{1.4}
\] In degree m its only possible summands have N<\allowbreak{}=p<\allowbreak{}=m. They form a finite sum of injectives. Column exactness in (1.3)\allowbreak{},\allowbreak{} and the finite-diagonal spectral comparison proved in the preceding note,\allowbreak{} show that epsilon:\allowbreak{}K\allowbreak{}->Tot(I)\allowbreak{} is a quasi-isomorphism. Each epsilon\^{}m is also monic:\allowbreak{} it is the monomorphism into I\textasciicircum{}\{m,\allowbreak{}0\} followed by the inclusion into that finite direct sum.

For a chain map f:\allowbreak{}K\allowbreak{}->L,\allowbreak{} use E\_\allowbreak{}q(f\textasciicircum{}p)\allowbreak{} on each summand. It commutes with both differentials by functoriality and naturality,\allowbreak{} hence with (1.4)\allowbreak{},\allowbreak{} and makes the augmentation square commute strictly. Identity and composition hold before passage to homotopy. This supplies the omitted verification at 7675–\allowbreak{}7677 and the functor inj to the actual arrow category. Different chosen lower bounds N merely certify the support:\allowbreak{} the formula (1.2)\allowbreak{} itself makes no choice of a bound for the object or for a map.

On homotopy classes the functor is the previously defined resolution functor j. If f and g are homotopic,\allowbreak{} epsilon\_\allowbreak{}L f and epsilon\_\allowbreak{}L g are homotopic. Precomposition with epsilon\_\allowbreak{}K is a bijection into the bounded-below injective target jL,\allowbreak{} so inj(f)\allowbreak{} and inj(g)\allowbreak{} are homotopic. The same argument,\allowbreak{} applied to inj(f\allowbreak{}+g)\allowbreak{} and inj(f)\allowbreak{}\allowbreak{}+inj(g)\allowbreak{},\allowbreak{} proves additivity of j. The latter two maps agree strictly after epsilon\_\allowbreak{}K. Their difference factors through jK/\allowbreak{}K,\allowbreak{} which is acyclic because epsilon\_\allowbreak{}K is a monomorphic quasi-isomorphism. A null homotopy of that quotient-to-jL map,\allowbreak{} pulled back along jK\allowbreak{}->jK/\allowbreak{}K,\allowbreak{} gives a homotopy h with h epsilon\_\allowbreak{}K\allowbreak{}=0. Thus this additivity comparison can be made relative to the original K; we do not assert a canonical selection of those homotopies.

There is a concrete reason not to assume strict additivity. On Q-vector spaces let J(V)\allowbreak{}\allowbreak{}=V direct-sum (V tensor\_\allowbreak{}Q V)\allowbreak{},\allowbreak{} eta(v)\allowbreak{}\allowbreak{}=(v,\allowbreak{}0)\allowbreak{},\allowbreak{} and J(f)\allowbreak{}\allowbreak{}=f direct-sum (f tensor f)\allowbreak{}. This is a zero-preserving functorial injective embedding,\allowbreak{} since every vector space is injective (extend linear maps from a basis of a subspace to a basis of the whole space)\allowbreak{}. For K\allowbreak{}=Q[0]\allowbreak{},\allowbreak{} (1.2)\allowbreak{} gives I\textasciicircum{}q\allowbreak{}=Q direct-sum Q for all q>\allowbreak{}=0 and d\textasciicircum{}q(x,\allowbreak{}y)\allowbreak{}\allowbreak{}=(y,\allowbreak{}0)\allowbreak{}. On a scalar lambda its degree-q map is diag(lambda\textasciicircum{}\{2\textasciicircum{}q\},\allowbreak{}lambda\textasciicircum{}\{2\textasciicircum{}\{q\allowbreak{}+1\}\})\allowbreak{}. For f\allowbreak{}=g\allowbreak{}=1,\allowbreak{} the additivity defect is therefore \[
 \Delta^q=
 \begin{pmatrix}2^{2^q}-2&0\\0&2^{2^{q+1}}-2\end{pmatrix}.
                                                               \tag{1.5}
\] In degree 0 it has nonzero second component 2. Put h\textasciicircum{}0\allowbreak{}=0 and for q>\allowbreak{}=1 set \[
 h^q(x,y)=(0,(2^{2^q}-2)x):I^q\longrightarrow I^{q-1}.
                                                               \tag{1.6}
\] Then d\^{}\{q-1\}h\^{}q contributes the first diagonal entry of (1.5)\allowbreak{},\allowbreak{} and h\textasciicircum{}\{q\allowbreak{}+1\}d\textasciicircum{}q the second. Their sum is exactly Delta,\allowbreak{} including q\allowbreak{}=0. The homotopy vanishes on K. This example proves both the failure at the chain level and the exact homotopy correction,\allowbreak{} with all coefficients retained.

RECON-052 reconciles P02-E0056 and DERIVED-048 with existing MC-STK-ERR-0859 (“remains”)\allowbreak{}. RECON-053 reconciles P02-E0057 and DERIVED-050 with the two MC-STK-ERR-0861 diagram edits to j(K\textasciicircum{}bullet)\allowbreak{} and j(L\textasciicircum{}bullet)\allowbreak{}. P02 also cites 7695,\allowbreak{} where the notation is already correct; that line is context and receives no extra edit.

\subsection{2. Large categories and the exact lift of a derived functor}\label{proofs-7589-8350-large-categories-and-the-exact-lift-of-a-derived-functor}

The construction (1.2)\allowbreak{} consists of specified functors,\allowbreak{} natural maps,\allowbreak{} and their cokernels,\allowbreak{} so it does not require selecting a resolution independently from a proper class of objects. The source\textquotesingle s applications use functorial injective embeddings for modules and sheaves. The current Injectives theorem at 1003–\allowbreak{}1014 is exactly the required assertion for sheaves on a ringed site. Its preceding displayed construction is J(F)\allowbreak{}\allowbreak{}=(FreeFlat(F\textasciicircum{}vee)\allowbreak{})\allowbreak{}\textasciicircum{}vee with the natural evaluation embedding. Its injectivity calculation identifies Hom(G,\allowbreak{}J(F)\allowbreak{})\allowbreak{} with Hom(G tensor FreeFlat(F\textasciicircum{}vee)\allowbreak{},\allowbreak{}J\_\allowbreak{}ab)\allowbreak{}; flatness of the first factor and injectivity of J\_\allowbreak{}ab make this exact. These are the particular construction and theorem cited by the receiver; the whole Injectives chapter was not re-reviewed here.

For presheaves,\allowbreak{} the topology consisting of isomorphism coverings has every presheaf as a sheaf:\allowbreak{} matching data on a single isomorphism pulls back uniquely along its inverse,\allowbreak{} and the same is true after any base change and composition. The sheaf theorem therefore yields the asserted presheaf embedding. Modules over a ring are sheaves of modules on the one-object,\allowbreak{} identity-morphism site with that ring; the same construction applies. These are the three completed examples. RECON-054 /\allowbreak{} P02-E0058 /\allowbreak{} DERIVED-049 retains MC-STK-ERR-0860\textquotesingle s deletion of the literal ``Add more here as needed.'' item. No new example is substituted into that placeholder.

Let F:\allowbreak{}A\allowbreak{}->B now be the additive functor of original 7778–\allowbreak{}7811,\allowbreak{} and (j,\allowbreak{}i)\allowbreak{} the resolution functor,\allowbreak{} with quasi-inverse j':\allowbreak{}D\textasciicircum{}\allowbreak{}+(A)\allowbreak{}\allowbreak{}->K\textasciicircum{}\allowbreak{}+(I)\allowbreak{}. Additive F preserves chain homotopies,\allowbreak{} finite direct sums,\allowbreak{} the cone differential including its minus sign,\allowbreak{} and shifts with their differential signs. It therefore gives an exact functor K\textasciicircum{}\allowbreak{}+(I)\allowbreak{}\allowbreak{}->K\textasciicircum{}\allowbreak{}+(B)\allowbreak{}. For K\allowbreak{}->jK the canonical derived comparison identifies RF(K)\allowbreak{} with F(jK)\allowbreak{}. Naturality follows from j(alpha)\allowbreak{}i\_\allowbreak{}K\allowbreak{}=i\_\allowbreak{}L alpha and the unique injective comparison; the same equation after localization handles inverse quasi-isomorphisms. Thus the diagram in the source 2-commutes,\allowbreak{} and the composite Fj\textquotesingle{} is an actual exact lift to K\textasciicircum{}\allowbreak{}+(B)\allowbreak{}. It is stronger data than its image in D\textasciicircum{}\allowbreak{}+(B)\allowbreak{},\allowbreak{} but no strict chain-level inverse to localization is asserted.

RECON-055 /\allowbreak{} P02-E0059 /\allowbreak{} DERIVED-051 is MC-STK-ERR-0862,\allowbreak{} the terminal period at 7780. RECON-056 /\allowbreak{} DERIVED-052 is MC-STK-ERR-0863,\allowbreak{} the defining order (j,\allowbreak{}i)\allowbreak{}.

\subsection{\texorpdfstring{3. Finite filtrations,\allowbreak{} split injectives,\allowbreak{} and the reversed pushout}{3. Finite  split  and the reversed pushout}}\label{proofs-7589-8350-finite-filtrations-split-injectives-and-the-reversed-pushout}

All filtrations in this section are decreasing,\allowbreak{} finite,\allowbreak{} exhaustive and separated. A map is strict when its image of F\^{}p equals its full image intersected with F\^{}p of the target. For a complex A\allowbreak{}->B\allowbreak{}->C,\allowbreak{} exactness of the associated graded complex implies exactness of each filtered piece and each quotient. Here is the finite induction used by current Homology 5184–\allowbreak{}5233. Remove the top common filtration level n. The three quotient objects have shorter filtrations; their graded complex is exact. Their quotient complexes are exact by induction. The complex of top pieces is exact as a graded piece. The short exact sequence of these three-term complexes then makes the original complex exact at its middle term by the cohomology long exact sequence. It also proves the quotient assertions,\allowbreak{} and the reversed finite induction proves the filtered piece assertions. Exactness on pieces and on the underlying objects gives f(F\^{}p A)\allowbreak{}\allowbreak{}=ker g intersect F\^{}p B for f:\allowbreak{}A\allowbreak{}->B,\allowbreak{} so f is strict. For g,\allowbreak{} if g(b)\allowbreak{} lies in F\^{}p C,\allowbreak{} exactness on quotients puts the class of b in the image of A/\allowbreak{}F\textasciicircum{}p A; subtracting that image leaves a representative in F\^{}p B with the same g-image. This proves strictness of g as an equality of image subobjects. These arguments can be expressed using the corresponding pullbacks and images in an arbitrary abelian category; no element-lifting surjectivity is imposed.

If each gr\^{}p I is injective,\allowbreak{} split 0\allowbreak{}->F\textasciicircum{}\{p\allowbreak{}+1\}I\allowbreak{}->F\textasciicircum{}p I\allowbreak{}->gr\textasciicircum{}p I\allowbreak{}->0 downwards from the top level. The kernel at each step is a finite sum of injectives,\allowbreak{} hence injective,\allowbreak{} so that sequence splits. We obtain I\allowbreak{}=direct-sum I\_\allowbreak{}p with F\^{}a I\allowbreak{}=direct-sum\_\allowbreak{}\{p>\allowbreak{}=a\}I\_\allowbreak{}p. Conversely such a decomposition has the required injective graded pieces. This proves 7851–\allowbreak{}7864 without requiring injectivity of a quotient to split an epimorphism.

Let u:\allowbreak{}A\allowbreak{}->B be a strict monomorphism and f:\allowbreak{}A\allowbreak{}->I filtered,\allowbreak{} with that finite splitting of I. The component f\_\allowbreak{}p to I\_\allowbreak{}p vanishes on F\textasciicircum{}\{p\allowbreak{}+1\}A,\allowbreak{} so factors through A/\allowbreak{}F\textasciicircum{}\{p\allowbreak{}+1\}A. Strictness of u makes that quotient a subobject of B/\allowbreak{}F\textasciicircum{}\{p\allowbreak{}+1\}B. Extend f\_\allowbreak{}p there by injectivity of I\_\allowbreak{}p. Pull it back to B,\allowbreak{} and assemble the finitely many components. For x in F\^{}a B every component of degree p\textless a vanishes,\allowbreak{} so this is a filtered extension g with gu\allowbreak{}=f. Taking A\allowbreak{}=I and f\allowbreak{}=1 gives the strict-monomorphism splitting lemma as well. For I\allowbreak{}=0 its splitting is the unique map; the source\textquotesingle s top-index induction is used only for I!\allowbreak{}=0.

The pushout used at 7906 is \[
 P=I\amalg_A B
   =\operatorname{coker}((f,-u):A\to I\oplus B),\qquad
 f':I\longrightarrow P.                                      \tag{3.1}
\] It is f\textquotesingle{} that is the pushout of u along f. Its underlying map is monic because u is. The quotient filtration is the image of F\^{}p I direct-sum F\^{}p B. If the image of x in I lies in F\^{}p P,\allowbreak{} the quotient relation says x-f(a)\allowbreak{} lies in F\^{}p I for some a with u(a)\allowbreak{} in F\^{}p B. Strictness of u gives a in F\^{}p A,\allowbreak{} so f(a)\allowbreak{} lies in F\^{}p I. Thus x lies in F\^{}p I and f\textquotesingle{} is strict. This is the specific implication of Homology 4988–\allowbreak{}5013 used here. A filtered retraction r:\allowbreak{}P\allowbreak{}->I exists by the preceding splitting lemma. If v:\allowbreak{}B\allowbreak{}->P is the other map,\allowbreak{} g\allowbreak{}=rv satisfies gu\allowbreak{}=rvu\allowbreak{}=rf'f\allowbreak{}=f.

\textbf{DERIVED-NEW-027.} Original 7906 calls f\textquotesingle{} the pushout ``of f by u''. Those roles are reversed:\allowbreak{} that would give the other map v,\allowbreak{} whose domain is B. The strict monomorphism follows from u,\allowbreak{} not from arbitrary f. Replace that phrase by ``of u by f''. For example,\allowbreak{} over Q-vector spaces with their single filtration step,\allowbreak{} take A\allowbreak{}=B\allowbreak{}=Q,\allowbreak{} u\allowbreak{}=1,\allowbreak{} I\allowbreak{}=0,\allowbreak{} f\allowbreak{}=0. Then P\allowbreak{}=0 and f':\allowbreak{}0\allowbreak{}->0 is monic,\allowbreak{} whereas the pushout of f along u is Q\allowbreak{}->0,\allowbreak{} which is not. The theorem and its displayed arrow are valid; the phrase identifies the wrong arrow operation.

For the embedding at 7913,\allowbreak{} choose a<\allowbreak{}=b with gr\^{}p A\allowbreak{}=0 outside [a,\allowbreak{}b]\allowbreak{} and monomorphisms w\_\allowbreak{}p:\allowbreak{}A/\allowbreak{}F\textasciicircum{}\{p\allowbreak{}+1\}A\allowbreak{}->I\_\allowbreak{}p. Define I\allowbreak{}=direct-sum\_\allowbreak{}\{a<\allowbreak{}=p<\allowbreak{}=b\}I\_\allowbreak{}p with the stated step filtration,\allowbreak{} and u\allowbreak{}=(w\_\allowbreak{}p pi\_\allowbreak{}p)\allowbreak{}\_\allowbreak{}p:\allowbreak{}A\allowbreak{}->I,\allowbreak{} where pi\_\allowbreak{}p is the actual quotient map. The b component is a monomorphism because F\textasciicircum{}\{b\allowbreak{}+1\}A\allowbreak{}=0,\allowbreak{} so u is monic. For a<r<\allowbreak{}=b\allowbreak{}+1,\allowbreak{} the condition u(x)\allowbreak{} in F\^{}r I forces its r-1 component to vanish; monicity of w\_\allowbreak{}\{r-1\} gives x in F\^{}r A. The converse uses pi\_\allowbreak{}p(F\textasciicircum{}r A)\allowbreak{}\allowbreak{}=0 for p\textless r. At r<\allowbreak{}=a both sides are the whole object; above b\allowbreak{}+1 both are zero by monicity. This proves strictness with the source\textquotesingle s quotients and direct-sum components,\allowbreak{} including the implicit quotient in its phrase ``sum of the maps u\_\allowbreak{}n”.

\subsection{\texorpdfstring{4. Filtered resolutions of objects,\allowbreak{} maps,\allowbreak{} and short exact sequences}{4. Filtered resolutions of   and short exact sequences}}\label{proofs-7589-8350-filtered-resolutions-of-objects-maps-and-short-exact-sequences}

Successively embed A,\allowbreak{} then its filtered cokernel,\allowbreak{} then the next cokernel into filtered injectives by section 3. Strictness of each embedding makes 0\allowbreak{}->Q\_\allowbreak{}n\allowbreak{}->I\textasciicircum{}\{n\allowbreak{}+1\}\allowbreak{}->Q\_\allowbreak{}\{n\allowbreak{}+1\}\allowbreak{}->0 exact on every graded piece. The differential is quotient followed by the next embedding,\allowbreak{} so its square is zero and its graded kernel is the preceding graded image. Thus A[0]\allowbreak{}\allowbreak{}->I is a filtered quasi-isomorphism,\allowbreak{} with zero negative terms. RECON-057 /\allowbreak{} P02-E0060 /\allowbreak{} DERIVED-054 retains both MC-STK-ERR-0865 corrections:\allowbreak{} the filtered injective objects belong to Fil\textasciicircum{}f(A)\allowbreak{}.

For f:\allowbreak{}A\allowbreak{}->B and the given resolutions A\allowbreak{}->I,\allowbreak{} B\allowbreak{}->J,\allowbreak{} their augmented complexes are graded-exact. Section 3 therefore makes the augmentations strict monic and every differential strict. Extend bf across A\allowbreak{}->I\textasciicircum{}0 to f\textasciicircum{}0:\allowbreak{}I\textasciicircum{}0\allowbreak{}->J\textasciicircum{}0. The map d\_\allowbreak{}J\textasciicircum{}0 f\^{}0 vanishes on A,\allowbreak{} so factors through coker(A\allowbreak{}->I\textasciicircum{}0)\allowbreak{}\allowbreak{}=coim d\_\allowbreak{}I\textasciicircum{}0 \allowbreak{}=im d\_\allowbreak{}I\textasciicircum{}0,\allowbreak{} with its induced filtration. Extend it from that strict subobject of I\^{}1 to f\textasciicircum{}1:\allowbreak{}I\textasciicircum{}1\allowbreak{}->J\textasciicircum{}1. At each subsequent step d\_\allowbreak{}J f\^{}n vanishes on the preceding image because the earlier chain equation and d\_\allowbreak{}J\textasciicircum{}2\allowbreak{}=0 give that equality. The same factorization and extension continue. This proves the strict commuting diagram at 7960–\allowbreak{}7974. RECON-058 /\allowbreak{} P02-E0061 /\allowbreak{} DERIVED-053 is MC-STK-ERR-0864,\allowbreak{} replacing the undefined C[0]\allowbreak{} by the supplied B[0]\allowbreak{}.

Now keep the source\textquotesingle s short exact sequence 0\allowbreak{}->A --alpha-\allowbreak{}-> B --beta-\allowbreak{}-> C\allowbreak{}->0,\allowbreak{} with its prescribed resolutions A\allowbreak{}->I and C\allowbreak{}->J. Extend a:\allowbreak{}A\allowbreak{}->I\textasciicircum{}0 to b:\allowbreak{}B\allowbreak{}->I\textasciicircum{}0. Since d\_\allowbreak{}I\textasciicircum{}0 b kills alpha,\allowbreak{} there is a unique bar b:\allowbreak{}C\allowbreak{}->I\textasciicircum{}1 with bar b beta\allowbreak{}=d\_\allowbreak{}I\textasciicircum{}0 b. Extend bar b across c:\allowbreak{}C\allowbreak{}->J\textasciicircum{}0 to delta\textasciicircum{}0:\allowbreak{}J\textasciicircum{}0\allowbreak{}->I\textasciicircum{}1. The equality d\_\allowbreak{}I\textasciicircum{}1 bar b\allowbreak{}=0 follows after composition with the epimorphism beta. Thus d\_\allowbreak{}I\textasciicircum{}1 delta\^{}0 kills c(C)\allowbreak{} and factors through im d\_\allowbreak{}J\textasciicircum{}0. Extend to delta\textasciicircum{}1:\allowbreak{}J\textasciicircum{}1\allowbreak{}->I\textasciicircum{}2. Recursively,\allowbreak{} the preceding identity and d\_\allowbreak{}I\textasciicircum{}2\allowbreak{}=0 give \[
 d_I^{r+1}\delta^r=\delta^{r+1}d_J^r
 \quad(r\ge0).                                               \tag{4.1}
\] Each extension is filtered,\allowbreak{} since im d\_\allowbreak{}J\textasciicircum{}r is a strict subobject. Set M\textasciicircum{}r\allowbreak{}=I\textasciicircum{}r direct-sum J\^{}r and \[
 d_M^r=
 \begin{pmatrix}
 d_I^r&(-1)^{r+1}\delta^r\\
 0&d_J^r
 \end{pmatrix}.                                              \tag{4.2}
\] Its off-diagonal square is (-1)\allowbreak{}\textasciicircum{}\{r\allowbreak{}+1\}d\_\allowbreak{}I\textasciicircum{}\{r\allowbreak{}+1\}delta\textasciicircum{}r \allowbreak{}+(-1)\allowbreak{}\textasciicircum{}\{r\allowbreak{}+2\}delta\textasciicircum{}\{r\allowbreak{}+1\}d\_\allowbreak{}J\textasciicircum{}r\allowbreak{}=0 by (4.1)\allowbreak{}. The augmentation B\allowbreak{}->M\textasciicircum{}0 is (b,\allowbreak{}c beta)\allowbreak{}. Its differential has first component d\_\allowbreak{}I\textasciicircum{}0 b-delta\^{}0 c beta \allowbreak{}=bar b beta-bar b beta\allowbreak{}=0 and second component zero. The maps from I and to J are the actual inclusion and projection. They give a termwise filtered split exact sequence,\allowbreak{} compatible with the original short exact sequence. Applying gr\^{}p gives a map of short exact sequences of complexes with outer quasi-isomorphisms. The cohomology exact sequences prove that the middle augmentation is a quasi-isomorphism on every graded piece.

RECON-059 /\allowbreak{} P02-E0062 /\allowbreak{} DERIVED-063 is the existing MC-STK-ERR-0874 chapter-prefix repair at 8044. The unprefixed label actually names a different Derived result; it is not an absent label as P02 says. The required result is the strictness lemma in Homology proved above. RECON-060 /\allowbreak{} P02-E0063 /\allowbreak{} DERIVED-055 is MC-STK-ERR-0866,\allowbreak{} the source B of (b,\allowbreak{}c beta)\allowbreak{}. P02\textquotesingle s earlier diagram loci 8049–\allowbreak{}8050 and 8060 are supporting context; the correction is at 8092 alone.

\subsection{5. Arbitrary bounded-below filtered complexes}\label{proofs-7589-8350-arbitrary-bounded-below-filtered-complexes}

Keep K\textasciicircum{}p\allowbreak{}=0 for p\textless N in its original degrees. For each p>\allowbreak{}=N form the source\textquotesingle s strict short exact sequence \[
 0\to\ker d_K^p\to K^p
       \to\operatorname{coim}d_K^p\to0.                      \tag{5.1}
\] The kernel has the induced filtration and the coimage the quotient filtration. No strictness of d\_\allowbreak{}K\textasciicircum{}p itself is assumed. Since d\_\allowbreak{}K\textasciicircum{}\{p\allowbreak{}+1\}d\_\allowbreak{}K\textasciicircum{}p\allowbreak{}=0,\allowbreak{} there is a filtered map u\_\allowbreak{}p:\allowbreak{}coim d\_\allowbreak{}K\textasciicircum{}p\allowbreak{}->ker d\_\allowbreak{}K\textasciicircum{}\{p\allowbreak{}+1\} through which the original differential factors. The image of F\^{}r K\^{}p lies in the induced F\^{}r of that kernel,\allowbreak{} so the factor is a filtered map with the quotient filtration on its domain.

Resolve the two outer objects in (5.1)\allowbreak{},\allowbreak{} lift u\_\allowbreak{}p by section 4,\allowbreak{} and apply its filtered horseshoe with those outer resolutions fixed. Call the resulting maps alpha\_\allowbreak{}p:\allowbreak{}I\_\allowbreak{}ker,\allowbreak{}p\allowbreak{}->I\_\allowbreak{}p and beta\_\allowbreak{}p:\allowbreak{}I\_\allowbreak{}p\allowbreak{}->I\_\allowbreak{}coim,\allowbreak{}p. The horizontal differential is exactly \[
 d_p=\alpha_{p+1}u_p^\bullet\beta_p.                          \tag{5.2}
\] Its composite with d\_\allowbreak{}\{p-1\} is zero because beta\_\allowbreak{}p alpha\_\allowbreak{}p\allowbreak{}=0. All factors in (5.2)\allowbreak{} are vertical chain maps,\allowbreak{} so it commutes with the vertical differentials. Every column is a filtered injective resolution of K\textasciicircum{}p,\allowbreak{} zero in negative vertical degree. Totalize with d\_\allowbreak{}1\allowbreak{}+(-1)\allowbreak{}\textasciicircum{}p d\_\allowbreak{}2. In degree m there are only the columns N<\allowbreak{}=p<\allowbreak{}=m. Every term is therefore filtered injective and has finite filtration. The augmentation K\textasciicircum{}m\allowbreak{}->I\_\allowbreak{}m\textasciicircum{}0 is strict monic,\allowbreak{} since its associated graded maps are monic at the bottom of their resolutions. The inclusion of I\_\allowbreak{}m\textasciicircum{}0 into the total finite direct sum is strict monic too. Their composite is the required component.

The associated graded functor is additive,\allowbreak{} hence commutes with these finite total sums and with the differential including (-1)\allowbreak{}\textasciicircum{}p. Applying the column comparison to each gr\^{}r proves that the actual augmentation K\allowbreak{}->Tot I is a filtered quasi-isomorphism. It is zero below N. This proves all three clauses of 8095–\allowbreak{}8103 without a change of the original degree coordinates. The source\textquotesingle s shifted proof is compatible:\allowbreak{} undoing its initial shift on both the source and constructed resolution multiplies the differentials by the prescribed shift sign. No strict commutation of a nonadditive embedding functor with shifts is asserted. We have instead specified the construction directly in the original p.

RECON-061 /\allowbreak{} DERIVED-056 is MC-STK-ERR-0867,\allowbreak{} inserting ``by'' in the assignment of the total complex notation.

\subsection{\texorpdfstring{6. Filtered null homotopies,\allowbreak{} comparison,\allowbreak{} and equivalence}{6. Filtered null   and equivalence}}\label{proofs-7589-8350-filtered-null-homotopies-comparison-and-equivalence}

Let C be filtered acyclic and I a complex of filtered injectives,\allowbreak{} zero below N. For a filtered chain map f:\allowbreak{}C\allowbreak{}->I set h\textasciicircum{}r\allowbreak{}=0 for r<\allowbreak{}=N. Suppose h through degree r has been constructed,\allowbreak{} and put \[
 e^r=f^r-d_I^{r-1}h^r.
\] The earlier homotopy equation and the chain equation give e\^{}r d\_\allowbreak{}C\textasciicircum{}\{r-1\}\allowbreak{}=0. Filtered exactness of C identifies coker d\_\allowbreak{}C\textasciicircum{}\{r-1\} with im d\_\allowbreak{}C\textasciicircum{}r,\allowbreak{} with its induced filtration. Thus e\^{}r factors as a filtered map on im d\_\allowbreak{}C\textasciicircum{}r. Extend it across its strict inclusion into C\textasciicircum{}\{r\allowbreak{}+1\} to h\textasciicircum{}\{r\allowbreak{}+1\}:\allowbreak{}C\textasciicircum{}\{r\allowbreak{}+1\}\allowbreak{}->I\textasciicircum{}r. Then \[
 f^r=d_I^{r-1}h^r+h^{r+1}d_C^r.                              \tag{6.1}
\] At each fixed degree only this step and the preceding step contribute,\allowbreak{} so the infinite upward induction defines an ordinary homotopy,\allowbreak{} not an infinite sum in A. This proves 8196–\allowbreak{}8238 and the filtered null-homotopy fact used in PART10 section 3.

For a filtered quasi-isomorphism alpha:\allowbreak{}K\allowbreak{}->L,\allowbreak{} its cone and shifted cone are filtered acyclic. Apply Hom(-,\allowbreak{}I)\allowbreak{} to the source\textquotesingle s triangle (K,\allowbreak{}L,\allowbreak{}C(alpha)\allowbreak{},\allowbreak{}alpha,\allowbreak{}i,\allowbreak{}-p)\allowbreak{}. The outer Hom groups vanish by (6.1)\allowbreak{},\allowbreak{} so precomposition is a bijection \[
 \operatorname{Hom}_{K(\operatorname{Fil}^f A)}(L,I)
 \xrightarrow{\sim}
 \operatorname{Hom}_{K(\operatorname{Fil}^f A)}(K,I).           \tag{6.2}
\] A derived roof gamma alpha\^{}\{-1\} is represented by the unique homotopy class beta with beta alpha\allowbreak{}=gamma. If a homotopy class becomes zero after localization,\allowbreak{} a denominator alpha makes its precomposition zero,\allowbreak{} and (6.2)\allowbreak{} makes it zero. This proves both surjectivity and injectivity of the Hom comparison with DF(A)\allowbreak{}. RECON-062 /\allowbreak{} P02-E0064 /\allowbreak{} DERIVED-057 retains the two MC-STK-ERR-0868 repairs of the ambient Hom category. The assertion is (6.2)\allowbreak{}; it does not identify filtered and unfiltered homotopy classes.

Finite sums and shifts of filtered injectives remain filtered injective,\allowbreak{} so K\textasciicircum{}\allowbreak{}+(I\textasciicircum{}f)\allowbreak{} is closed under the actual cone triangles. Its functor to DF\textasciicircum{}\allowbreak{}+(A)\allowbreak{} is exact,\allowbreak{} and (6.2)\allowbreak{} makes it fully faithful. An object of DF\textasciicircum{}\allowbreak{}+(A)\allowbreak{} has a representative with gr cohomology zero below some N. The already proved bounded-representative lemma at original 4385–\allowbreak{}4431 replaces that representative by a filtered quasi-isomorphic complex zero below N:\allowbreak{} strictness of d\^{}\{N-1\} justifies the quotient K\textasciicircum{}N/\allowbreak{}im d\^{}\{N-1\} at its bottom degree. Applying section 5 then supplies a bounded-below filtered injective representative. Thus the functor is essentially surjective. Both gr\^{}p and forgetting the filtration send a filtered injective to an injective object,\allowbreak{} by its finite splitting. They preserve shifts,\allowbreak{} cones and filtered quasi-isomorphisms. The two squares in 8315–\allowbreak{}8335 therefore commute with their displayed functors.

\subsection{\texorpdfstring{7. DERIVED-UNDER-013:\allowbreak{} degree bounds and functorial fixed filtration windows}{7.  degree bounds and functorial fixed filtration windows}}\label{proofs-7589-8350-derived-under-013-degree-bounds-and-functorial-fixed-filtration-windows}

The construction in section 5 can retain more than just finiteness. Suppose K\textasciicircum{}p,\allowbreak{} p>\allowbreak{}=N,\allowbreak{} has filtration in the interval [a\_\allowbreak{}p,\allowbreak{}b\_\allowbreak{}p]\allowbreak{},\allowbreak{} meaning F\textasciicircum{}\{a\_\allowbreak{}p\}K\textasciicircum{}p\allowbreak{}=K\textasciicircum{}p and F\textasciicircum{}\{b\_\allowbreak{}p\allowbreak{}+1\}K\textasciicircum{}p\allowbreak{}=0. Its kernel and coimage in (5.1)\allowbreak{} have the same permitted bounds. Section 3\textquotesingle s embedding uses precisely the indices a\_\allowbreak{}p<\allowbreak{}=r<\allowbreak{}=b\_\allowbreak{}p,\allowbreak{} and its cokernel retains that interval. Iterating gives both outer resolutions in that same interval. Their horseshoe middle is their direct sum in each degree,\allowbreak{} so retains it too. Consequently the resulting total resolution has in degree m>\allowbreak{}=N the explicit interval \[
 \left[\min_{N\le p\le m}a_p,\
       \max_{N\le p\le m}b_p\right].                         \tag{7.1}
\] No term with p\textgreater m contributes. This is a finite interval for each m,\allowbreak{} even if the filtration widths of K have no common global bound. All terms below N remain zero.

If A has a functorial injective embedding,\allowbreak{} fix a single interval [a,\allowbreak{}b]\allowbreak{}. On the category of filtered objects supported in that interval define the functor \[
 \mathcal E(A)=
 \bigoplus_{r=a}^{b}J_{\rm red}(A/F^{r+1}A),\qquad
 F^p\mathcal E(A)=
 \bigoplus_{\substack{a\le r\le b\\r\ge p}}
              J_{\rm red}(A/F^{r+1}A).                       \tag{7.2}
\] The natural map is the tuple of quotient maps followed by eta\_\allowbreak{}red. The proof at the end of section 3 shows it is strict monic. Each graded piece is injective,\allowbreak{} so the target is filtered injective. Quotient maps of filtered morphisms give the functor on maps. It preserves zero because J\_\allowbreak{}red does,\allowbreak{} and the filtered cokernel of its natural embedding stays in [a,\allowbreak{}b]\allowbreak{}.

Repeat (1.2)\allowbreak{} in this filtered category,\allowbreak{} using (7.2)\allowbreak{} and these cokernels. Every augmented column is exact on each graded piece,\allowbreak{} because every embedding is strict. For a bounded-below complex all of whose terms have this fixed filtration interval,\allowbreak{} the horizontal maps are the images of its actual differentials under these zero-preserving functors. Their squares are zero,\allowbreak{} and naturality commutes the two differentials. The total construction (1.4)\allowbreak{} gives a functorial strict monomorphic filtered quasi-isomorphism,\allowbreak{} retains the original lower term bound N,\allowbreak{} and retains [a,\allowbreak{}b]\allowbreak{} in every degree. It uses the specified functors and no selection of resolutions over a proper class of objects.

The two constructions have different assertions:\allowbreak{} (7.1)\allowbreak{} covers arbitrary finite filtrations,\allowbreak{} with the source\textquotesingle s choices; (7.2)\allowbreak{} is functorial on a fixed filtration window. A window chosen separately for each object has not been shown to yield a functor on all finite filtered objects. We do not make that assertion.

The resulting homotopy functor is additive and exact by section 6. Its strict chain functor need not be additive,\allowbreak{} as (1.5)\allowbreak{} proves already with a single filtration step. The additivity defect has a filtered homotopy relative to the source:\allowbreak{} its factor through the cokernel of the strict augmentation has filtered acyclic domain,\allowbreak{} so (6.1)\allowbreak{} gives that homotopy. This adds an exact comparison to the stronger support and functorial assertions; no higher coherence of chosen homotopies is claimed.

\subsection{\texorpdfstring{8. DERIVED-UNDER-014:\allowbreak{} both prescribed resolutions for filtered complexes}{8.  both prescribed resolutions for filtered complexes}}\label{proofs-7589-8350-derived-under-014-both-prescribed-resolutions-for-filtered-complexes}

The source\textquotesingle s horseshoe in section 4 starts with objects in degree zero. It extends to any degreewise strict short exact sequence of bounded-below filtered complexes \[
 0\to A\xrightarrow{u}B\xrightarrow{v}C\to0,                 \tag{8.1}
\] with any two given filtered quasi-isomorphisms f\_\allowbreak{}1:\allowbreak{}A\allowbreak{}->I\_\allowbreak{}1 and f\_\allowbreak{}3:\allowbreak{}C\allowbreak{}->I\_\allowbreak{}3 into bounded-below filtered injective complexes. Neither given resolution map must be a termwise monomorphism.

Form the filtered-complex pushout \[
 E=\operatorname{coker}
       ((f_1,-u):A\to I_1\oplus B).                          \tag{8.2}
\] Section 3 makes 0\allowbreak{}->I\_\allowbreak{}1\allowbreak{}->E\allowbreak{}->C\allowbreak{}->0 degreewise strict exact. Apply gr\^{}p to the map from (8.1)\allowbreak{} into this sequence. Its outer maps are gr\^{}p f\_\allowbreak{}1 and the identity,\allowbreak{} so B\allowbreak{}->E is a filtered quasi-isomorphism. Split I\_\allowbreak{}1\textasciicircum{}n\allowbreak{}->E\textasciicircum{}n in the filtered category. This identifies \[
 E^n=I_1^n\oplus C^n,\qquad
 d_E^n=\begin{pmatrix}d_{I_1}^n&\delta^n\\0&d_C^n\end{pmatrix}.
                                                               \tag{8.3}
\] Every entry is filtered. The squared differential gives d\_\allowbreak{}\{I\_\allowbreak{}1\}delta\allowbreak{}+delta d\_\allowbreak{}C\allowbreak{}=0,\allowbreak{} so delta:\allowbreak{}C\allowbreak{}->I\_\allowbreak{}1[1]\allowbreak{} is a filtered chain map with the original positive off-diagonal entry. Apply (6.2)\allowbreak{} to f\_\allowbreak{}3 and target I\_\allowbreak{}1[1]\allowbreak{}. Choose its preimage delta':\allowbreak{}I\_\allowbreak{}3\allowbreak{}->I\_\allowbreak{}1[1]\allowbreak{} and a filtered homotopy h\textasciicircum{}n:\allowbreak{}C\textasciicircum{}n\allowbreak{}->I\_\allowbreak{}1\textasciicircum{}n with \[
 (\delta')^n f_3^n-\delta^n
      =-d_{I_1}^n h^n+h^{n+1}d_C^n.                         \tag{8.4}
\] The minus sign comes from d\_\allowbreak{}\{I\_\allowbreak{}1[1]\allowbreak{}\}. Set \[
 M^n=I_1^n\oplus I_3^n,\quad
 d_M^n=\begin{pmatrix}d_{I_1}^n&(\delta')^n\\0&d_{I_3}^n\end{pmatrix},
 \quad
 e^n=
 \begin{pmatrix}1&h^n\\0&f_3^n\end{pmatrix}:E^n\to M^n.        \tag{8.5}
\] The chain equation for delta\textquotesingle{} proves d\_\allowbreak{}M\textasciicircum{}2\allowbreak{}=0. The two off-diagonal entries in d\_\allowbreak{}M e and e d\_\allowbreak{}E are respectively d\_\allowbreak{}\{I\_\allowbreak{}1\}h\allowbreak{}+delta'f\_\allowbreak{}3 and delta\allowbreak{}+h d\_\allowbreak{}C; (8.4)\allowbreak{} makes them equal. Thus e is a filtered chain map. Its outer maps in the split sequences are 1 and f\_\allowbreak{}3,\allowbreak{} so gr\^{}p e is a quasi-isomorphism for every p by the cohomology exact sequences. Composing B\allowbreak{}->E\allowbreak{}->M gives the middle filtered resolution,\allowbreak{} compatible with both prescribed outer ones. Its terms are filtered injective,\allowbreak{} and it retains any common lower term bound and common filtration interval of I\_\allowbreak{}1,\allowbreak{}I\_\allowbreak{}3. No extra global enough-injectives hypothesis is used once these resolutions are supplied.

This is the exact filtered counterpart of DERIVED-UNDER-010,\allowbreak{} including the homotopy matrix h. Replacing e by the diagonal matrix would be unjustified unless delta'f\_\allowbreak{}3\allowbreak{}=delta strictly. The source\textquotesingle s degree-zero construction already obtains that strict equation in its own signed form (4.1)\allowbreak{}–\allowbreak{}(4.2)\allowbreak{}; it is retained,\allowbreak{} not replaced.

\subsection{9. Receiving calculations and limits}\label{proofs-7589-8350-receiving-calculations-and-limits}

Current Cohomology 104–\allowbreak{}145 and Sites Cohomology 145–\allowbreak{}182 construct RF using resolution functors on their actual module categories. Sections 1–\allowbreak{}2 justify the strict chain construction,\allowbreak{} its descent to additive homotopy functors,\allowbreak{} and its exact lift F j\textquotesingle{} in those categories. Their F is left exact and therefore additive; no extra condition is required. The typo on the next Sites Cohomology line 183 is a separate receiving-source observation,\allowbreak{} queued for that chapter\textquotesingle s report reconciliation. It changes none of the mathematical maps reviewed here.

Current Cohomology 7310–\allowbreak{}7339 uses a bounded-below filtered injective resolution with a finite filtration on every term. Formula (7.1)\allowbreak{} proves exactly that finiteness even when the original filtration intervals vary with degree. If the original terms share [a,\allowbreak{}b]\allowbreak{},\allowbreak{} (7.2)\allowbreak{} gives a functorial choice with that same interval. The filtered complex of global sections uses F\^{}p Gamma(X,\allowbreak{}I\textasciicircum{}m)\allowbreak{} \allowbreak{}=Gamma(X,\allowbreak{}F\textasciicircum{}p I\textasciicircum{}m)\allowbreak{}; the split filtration makes these literal subobjects and keeps the same bounds. Thus its associated graded complexes and their cohomology use the same columns and signs proved above.

PART10 section 3 used filtered comparison to prove Cartan–\allowbreak{}Eilenberg functoriality and the decalage formula. Section 6 now supplies the full source-ordered review of that comparison and its category correction. Section 5 supplies the general filtered resolution existence needed in that forward passage. The proof-specific pushout correction propagates through the extension,\allowbreak{} map,\allowbreak{} horseshoe and total resolution constructions in sections 3–\allowbreak{}5. No theorem statement is weakened or silently strengthened by that correction.

The subsequent large-category remark and general filtered derived functor discussion were read through original 8420 but are next in the semantic review. Original 8407 already announces the additive filtered derived-functor treatment. That is relevant prior source evidence for PART10\textquotesingle s additive consequence; it reinforces the absence of a novelty claim. No complete review of the unread continuation or of the entire Injectives chapter is asserted here.
